\chapter{Discussion}
\label{ch:discussion}

\section{When target membership helps}

A target-linked representation can preserve which process outcome
belongs to which object as events from unrelated processes arrive.
The held-configuration comparison and the correct-versus-exchanged
membership control support that interpretation. The four-target
setting provides evidence that the model can exploit the relation
with unfamiliar target names. The parser still depends on observable
paths and a reliable START--END association. It can therefore provide
little extra information on prefixes with no recoverable target.

The timing contract has a separate consequence. Including a result
because its process has already exited can make the evidence look
stronger than it would be at a real prediction boundary. Preserving
the receipt boundary during training and evaluation exposes this
ambiguity. For an event with a received START and a delayed END,
the model must use the available prefix. An operational deployment
would also need a defined trigger for the forecast. The instrumented
source request supplies the trigger in these experiments.

\section{Interpreting improvements under benign activity}

The effect of benign interleaving varies across workload families.
In the B+C repeats, the within-architecture control supports
using the history state available at each event. In A+C, the
fixed two-model average reduced NLL on newly collected repetitions
under the specified training seeds. On the same runs, the wider
single-history and BiTA-inspired controls narrowed this difference
when all models were retrained. The comparisons therefore suggest
separate conclusions about a useful fixed predictor and the
repeatability of the architectural advantage.

A practical system would also account for inference cost.
For the fixed A+C system, the reported CPU measurement for a
batch of 32 prepared rows was a median of 9.32 milliseconds for the
baseline and 30.91 milliseconds for the two-model combination.
Those timings concern batched inference in the measured environment.
The operational cost would also include log receipt, preprocessing
and model loading where applicable.

\section{Threats to validity}

\textbf{Population.}
The study uses controlled Windows executions of known source
procedures and specified benign workloads. Separate repetitions
reduce sensitivity to one recording, while unfamiliar paths and
target counts change particular aspects of the input. Transfer to
unseen procedures, hosts in other organizations, and arbitrary
benign work remains an empirical question.

\textbf{Evidence availability.}
Command-line visibility and reliable process-event linkage depend
on the auditing setup and the received data. The application-specific
exit status is evidence about the corresponding process. The
source-designated label describes the executed behavior, whose
category can remain ambiguous from indirect audit evidence.
Windows logs may also carry sensitive command arguments in
deployment~\cite{ms4688}, making handling and retention of the
input a consideration.

\textbf{Repeated selection.}
Several model decisions were made during iterative development.
Separately marked frozen evaluations estimate performance on new
runs from known conditions. Post-hoc comparisons reuse some of
these same runs. The reported intervals describe the sampled
repetitions, configurations or training seeds stated with each
comparison. The uncertainty introduced by iterative model
selection requires an additional holdout to estimate.

\textbf{Verification.}
The quantitative tables are drawn from completed experiment reports.
Those reports describe saved-prediction checks and, for the later
evaluations, exact matching of the baseline output on target-free
prefixes. Independent replication from event manifests and checkpoints
remains open. A future field evaluation should specify the trigger
policy, define its label acquisition procedure, and retain the
receipt-time evidence contract.
